%% FullCourse_10Lecs -- Lecture 9, Part 9.1
%% Assembled by tools/lec09_assemble.py from reorg_manifest.yaml: Phase A of
%% Lec09_Reorg_Plan_2026-09-12.md as amended by Lec09_Reorg_Plan_Review_2026-09-12.md.
%% Each "% ---- <ID> · <TAG> · TIER" marker names a frame's source deck and line;
%% keep the markers until Phase B is done (lec09_assemble.py check reads them).
%% Owns: the example ladder (capability rungs); the principal-agent lens (Lec10 T9 owns the evidence)
%% Owns: chat vs agent (where the loop runs), the agentic loop, four features, the spectrum
%% Owns: orchestration-pattern names, multi-agent cost-benefit, measured task horizons (concept only)
%% Owns: the four research patterns (Lec10 T1 cites the names)
%% Defers: JSON tool call, session-as-list -> T2a; Git -> T2b; seven tests -> T5d
%% Defers: MCP is named once in the timeline and explained only in TA; model-specific numbers -> TA

%% Shared preamble for the FullCourse_10Lecs topic decks.
%%
%% Each topic .tex \input's this file, then defines its own \title, \date
%% override (if any), \graphicspath, and \begin{document}.
%%
%% Reconstructed verbatim from the inlined copy preserved in
%% Lec10_Case_Studies/Lec10_T8_Case6_DeepRamsey.tex (lines 23-190), which is
%% explicitly annotated there as "from ../shared_preamble.tex, copied verbatim".
%% Base preamble matches MiniCourse_8hr/Slides/shared_preamble.tex; this version
%% additionally provides \sectiondivider, the conceptbox/cautionbox/examplebox/
%% takeawaybox tcolorboxes, and \compacttables used across the split decks.

\documentclass[10pt,english,aspectratio=169]{beamer}

\usepackage{ae,aecompl}
\usepackage[T1]{fontenc}
\usepackage{textcomp}
\usepackage[utf8]{inputenc}
\usepackage{babel}
\usepackage{datetime2}

\setcounter{secnumdepth}{3}
\setcounter{tocdepth}{3}

\usepackage{amsmath, amssymb, amsthm, graphicx}
\usepackage{booktabs, multirow}
\usepackage{tikz}
\usepackage{pgfplots}
\pgfplotsset{compat=1.16}
\usepackage[most]{tcolorbox}
\tcbuselibrary{skins, breakable, theorems}
\usepackage{listings}
\usepackage{xcolor}

\lstset{
  basicstyle=\ttfamily\small,
  keywordstyle=\color{blue!80!black}\bfseries,
  commentstyle=\color{green!50!black}\itshape,
  stringstyle=\color{orange!80!black},
  backgroundcolor=\color{gray!8},
  frame=single,
  rulecolor=\color{gray!40},
  breaklines=true,
  showstringspaces=false,
  numbers=none,
}

\lstdefinestyle{pythonstyle}{
    language=Python,
    basicstyle=\ttfamily\footnotesize,
    keywordstyle=\color{blue!80!black}\bfseries,
    commentstyle=\color{green!50!black}\itshape,
    stringstyle=\color{orange!80!black},
    frame=single,
    breaklines=true,
    showstringspaces=false,
    numbers=left,
    numbersep=5pt,
    numberstyle=\tiny\color{gray},
    backgroundcolor=\color{gray!8},
    rulecolor=\color{gray!40}
}

\ifx\hypersetup\undefined
  \AtBeginDocument{%
    \hypersetup{bookmarks=false,breaklinks=true,pdfborder={0 0 0},colorlinks=true,
      linkcolor=DarkText,urlcolor=RedTitle}%
  }
\else
  \hypersetup{bookmarks=false,breaklinks=true,pdfborder={0 0 0},colorlinks=true,
    linkcolor=DarkText,urlcolor=RedTitle}%
\fi

\makeatletter
\newcommand\makebeamertitle{%
  \begin{frame}[plain]
    \vspace*{1.0cm}
    \centering
    {\usebeamerfont{title}\usebeamercolor[fg]{title}\inserttitle\par}
    \vspace{1.0cm}
    {\usebeamerfont{author}\usebeamercolor[fg]{author}\insertauthor\par}
    \vspace{0.2cm}
    {\usebeamerfont{date}\usebeamercolor[fg]{date}\insertdate\par}
  \end{frame}}%
\AtBeginDocument{%
  \let\origtableofcontents=\tableofcontents
  \def\tableofcontents{\@ifnextchar[{\origtableofcontents}{\gobbletableofcontents}}
  \def\gobbletableofcontents#1{\origtableofcontents}%
}

\usetheme{metropolis}
\usefonttheme{professionalfonts}

\definecolor{LightGray}{RGB}{230,230,230}
\definecolor{RedTitle}{RGB}{0,0,200}
\definecolor{VeryLightGray}{RGB}{245,245,245}
\definecolor{DarkText}{RGB}{0,0,0}

\setbeamercolor{background canvas}{bg=white}
\setbeamercolor{title}{fg=RedTitle}
\setbeamercolor{author}{fg=DarkText}
\setbeamercolor{date}{fg=DarkText}
\setbeamercolor{frametitle}{bg=VeryLightGray,fg=DarkText}
\setbeamercolor{normal text}{fg=DarkText}
\setbeamerfont{title}{size=\large}

\setbeamersize{text margin left=5mm, text margin right=5mm}
\addtobeamertemplate{frametitle}{}{\vspace{-0.5em}}
\newcommand{\reducevspace}{\vspace{-0.5cm}}

% Shared section divider and box styles for split topic decks.
\newcommand{\sectiondivider}[2]{%
\begin{frame}[plain,noframenumbering]
\begin{center}
\vfill
{\Large\textbf{#1}\par}
\vspace{0.35cm}
{\normalsize #2\par}
\vfill
\end{center}
\end{frame}
}

\newtcolorbox{conceptbox}[1][]{
  colback=blue!3,
  colframe=RedTitle!55,
  boxrule=0.35mm,
  arc=1mm,
  left=2mm,
  right=2mm,
  top=1mm,
  bottom=1mm,
  #1
}
\newtcolorbox{cautionbox}[1][]{
  colback=red!3,
  colframe=red!50!black,
  boxrule=0.35mm,
  arc=1mm,
  left=2mm,
  right=2mm,
  top=1mm,
  bottom=1mm,
  #1
}
\newtcolorbox{examplebox}[1][]{
  colback=green!3,
  colframe=green!45!black,
  boxrule=0.35mm,
  arc=1mm,
  left=2mm,
  right=2mm,
  top=1mm,
  bottom=1mm,
  #1
}
\newtcolorbox{takeawaybox}[1][]{
  colback=VeryLightGray,
  colframe=RedTitle!55,
  boxrule=0.35mm,
  arc=1mm,
  left=2mm,
  right=2mm,
  top=1mm,
  bottom=1mm,
  #1
}
\newcommand{\compacttables}{\renewcommand{\arraystretch}{1.15}}

\setbeamertemplate{navigation symbols}{}
\setbeamertemplate{footline}{%
  \hfill%
  \usebeamercolor[fg]{page number in head/foot}%
  \usebeamerfont{page number in head/foot}%
  \insertframenumber\,/\,\inserttotalframenumber\kern1em\vskip2pt%
}

\makeatother

\author{Zhigang Feng}
% "date with time": compile timestamp so each build is identifiable.
\date{\today\ \DTMcurrenttime}


% Suppress redundant auto-generated metropolis section page; \sectiondivider frames are the dividers we keep.
\metroset{sectionpage=none}

\graphicspath{{./pic/}{./figures/}{../}}

\usetikzlibrary{positioning,arrows.meta,shapes,calc,fit,backgrounds}

%% Lec09 shared MAP slide, v2 (September 2026 reorganization): five parts,
%% eleven decks.  v1 (the July "five acts" map) is archived as
%% _archive/five_acts_2026-07/lec09_map_v1.tex.
%%
%% Input this file AFTER ../shared_preamble.tex (needs RedTitle, VeryLightGray,
%% DarkText) and call \LecNineMap{part}{sub} inside a frame:
%%   part in {1..5} highlights that part ("you are here") and, when sub names a
%%   sub-deck letter, that deck's chip -- \LecNineMap{2}{b} is T2b, \LecNineMap{4}{}
%%   is T4;  part = 0 renders the closing reprise with every part complete (the
%%   "Your Job Description" frame in T5d).
%% Each part carries its student question and its economic twin (the delegation-
%% economics thread of Lec09_Reorg_Plan_Review_2026-09-12.md, A3).
%% Filename deliberately does not match the build glob Lec*_T*.tex, so the build
%% script never tries to compile it standalone.

% Highlight chip #1 when the requested deck key equals #2 (e.g. 2b).
\newcommand{\lecmapchip}[2]{%
  \edef\lecmaptmp{#2}%
  \ifx\lecmapkey\lecmaptmp\tikzset{#1/.style={lecchipcur}}\fi}

\newcommand{\LecNineMap}[2]{%
% TikZ option lists are comma-split before TeX conditionals run, so every
% \ifnum/\ifx is resolved here, into named styles, before the picture starts.
\edef\lecmapkey{#1#2}%
\tikzset{lecparta/.style={}, lecpartb/.style={}, lecpartc/.style={},
         lecpartd/.style={}, lecparte/.style={},
         chipTone/.style={}, chipTtwoa/.style={}, chipTtwob/.style={},
         chipTthreea/.style={}, chipTthreeb/.style={}, chipTfour/.style={},
         chipTfivea/.style={}, chipTfiveb/.style={}, chipTfivec/.style={},
         chipTfived/.style={}}%
\ifnum#1=0
  \tikzset{lecparta/.style={lecpartdone}, lecpartb/.style={lecpartdone},
           lecpartc/.style={lecpartdone}, lecpartd/.style={lecpartdone},
           lecparte/.style={lecpartdone}, lecchip/.style={lecchipbase, lecchipdone}}%
\else
  \tikzset{lecchip/.style={lecchipbase}}%
  \ifnum#1=1 \tikzset{lecparta/.style={lecpartcur}}\fi
  \ifnum#1=2 \tikzset{lecpartb/.style={lecpartcur}}\fi
  \ifnum#1=3 \tikzset{lecpartc/.style={lecpartcur}}\fi
  \ifnum#1=4 \tikzset{lecpartd/.style={lecpartcur}}\fi
  \ifnum#1=5 \tikzset{lecparte/.style={lecpartcur}}\fi
  \lecmapchip{chipTone}{1}\lecmapchip{chipTtwoa}{2a}\lecmapchip{chipTtwob}{2b}%
  \lecmapchip{chipTthreea}{3a}\lecmapchip{chipTthreeb}{3b}\lecmapchip{chipTfour}{4}%
  \lecmapchip{chipTfivea}{5a}\lecmapchip{chipTfiveb}{5b}%
  \lecmapchip{chipTfivec}{5c}\lecmapchip{chipTfived}{5d}%
\fi
\begin{center}
\begin{tikzpicture}[
    part/.style={rectangle, rounded corners, draw=black!60, fill=VeryLightGray,
        minimum width=2.62cm, minimum height=0.72cm, text width=2.5cm,
        align=center, font=\scriptsize, text=DarkText},
    lecpartcur/.style={draw=RedTitle, very thick, fill=orange!20},
    lecpartdone/.style={draw=green!45!black, thick, fill=green!10},
    lecchipbase/.style={rectangle, rounded corners=1pt, draw=black!30, fill=white,
        inner xsep=1.5pt, inner ysep=1pt, font=\tiny, text=black!70},
    lecchipcur/.style={draw=RedTitle, fill=RedTitle!12, text=RedTitle, font=\tiny\bfseries},
    lecchipdone/.style={draw=green!45!black, fill=green!8},
    q/.style={font=\tiny\itshape, text width=2.7cm, align=center, text=black!75},
    twin/.style={font=\tiny, text width=2.7cm, align=center, text=blue!40!black},
    sat/.style={rectangle, rounded corners, draw=black!45, dashed, fill=white,
        text width=5.6cm, align=center, font=\tiny, text=black!70, minimum height=0.42cm},
    barlab/.style={font=\tiny\bfseries, text=black!60},
    arr/.style={-{Stealth[length=2mm]}, thick, gray!60}
]
    % --- five part boxes ---
    \node[part, lecparta] (p1) at (0,0)     {\textbf{9.1 SEE}};
    \node[part, lecpartb] (p2) at (2.85,0)  {\textbf{9.2 UNDER THE HOOD}};
    \node[part, lecpartc] (p3) at (5.7,0)   {\textbf{9.3 BUILD}};
    \node[part, lecpartd] (p4) at (8.55,0)  {\textbf{9.4 DESIGN \& VALIDATE}};
    \node[part, lecparte] (p5) at (11.4,0)  {\textbf{9.5 RUN \& GOVERN}};
    \draw[arr] (p1) -- (p2);
    \draw[arr] (p2) -- (p3);
    \draw[arr] (p3) -- (p4);
    \draw[arr] (p4) -- (p5);
    % --- sub-deck chips ---
    \node[lecchip, chipTone]    at (0,-0.6)     {T1};
    \node[lecchip, chipTtwoa]   at (2.55,-0.6)  {T2a};
    \node[lecchip, chipTtwob]   at (3.15,-0.6)  {T2b};
    \node[lecchip, chipTthreea] at (5.4,-0.6)   {T3a};
    \node[lecchip, chipTthreeb] at (6.0,-0.6)   {T3b};
    \node[lecchip, chipTfour]   at (8.55,-0.6)  {T4};
    \node[lecchip, chipTfivea]  at (10.5,-0.6)  {T5a};
    \node[lecchip, chipTfiveb]  at (11.1,-0.6)  {T5b};
    \node[lecchip, chipTfivec]  at (11.7,-0.6)  {T5c};
    \node[lecchip, chipTfived]  at (12.3,-0.6)  {T5d};
    % --- the student question each part answers ---
    \node[q] at (0,-1.08)     {what changed,\\ and why care?};
    \node[q] at (2.85,-1.08)  {how does it run?\\ how do I start?};
    \node[q] at (5.7,-1.08)   {how do I build one,\\ and call it?};
    \node[q] at (8.55,-1.08)  {can I design one\\ and prove it works?};
    \node[q] at (11.4,-1.08)  {run a project, catch\\ mistakes, stay safe};
    % --- its economic twin ---
    \node[twin] at (0,-1.62)     {generation got cheap;\\ verification is scarce};
    \node[twin] at (2.85,-1.62)  {what am I paying for?\\ tokens, scarce context};
    \node[twin] at (5.7,-1.62)   {dividing the labour:\\ specialists, boundaries};
    \node[twin] at (8.55,-1.62)  {write the contract,\\ audit the work};
    \node[twin] at (11.4,-1.62)  {hidden action: gates,\\ monitoring, priors};
    % --- "you are here" marker above the current part ---
    \ifnum#1=1 \node[font=\scriptsize\bfseries, text=RedTitle] at (0,0.62)
        {you are here\ $\blacktriangledown$};\fi
    \ifnum#1=2 \node[font=\scriptsize\bfseries, text=RedTitle] at (2.85,0.62)
        {you are here\ $\blacktriangledown$};\fi
    \ifnum#1=3 \node[font=\scriptsize\bfseries, text=RedTitle] at (5.7,0.62)
        {you are here\ $\blacktriangledown$};\fi
    \ifnum#1=4 \node[font=\scriptsize\bfseries, text=RedTitle] at (8.55,0.62)
        {you are here\ $\blacktriangledown$};\fi
    \ifnum#1=5 \node[font=\scriptsize\bfseries, text=RedTitle] at (11.4,0.62)
        {you are here\ $\blacktriangledown$};\fi
    \ifnum#1=0 \node[font=\scriptsize\bfseries, text=green!45!black] at (5.7,0.62)
        {all five parts complete};\fi
    % --- bottom band: the three questions of the lecture ---
    \fill[blue!10]   (-1.3,-2.37) rectangle (4.15,-2.07);
    \fill[green!10]  (4.4,-2.37)  rectangle (9.85,-2.07);
    \fill[orange!15] (10.1,-2.37) rectangle (12.7,-2.07);
    \node[barlab] at (1.425,-2.22)  {HOW IT WORKS};
    \node[barlab] at (7.125,-2.22)  {HOW TO BUILD \& USE IT};
    \node[barlab] at (11.4,-2.22)   {YOUR ROLE};
    % --- dashed satellites ---
    \node[sat] at (1.9,-2.8)
        {\textbf{Appendix TA} --- the dated landscape and setup details, read on demand};
    \node[sat] at (9.9,-2.8)
        {\textbf{Lab} Parts A--B, Design Lab scenarios $\to$ \textbf{Lec10} cases (same morning)};
\end{tikzpicture}
\end{center}
}


\title[AI for Econ Research]{AI for Economic Research: Dynamic Models, Language, and Agents\\
Lecture 9.1: The Agent Opportunity --- From Chatbox to Agent Swarms}

\begin{document}

\makebeamertitle

% ---- NEW · TIER: CORE
% TODO(Phase B): Cold open before the map (archive_pre_split/Lec09_Agentic_AI.tex:2022-2121, "Case Study: Empty Folder -> Figure").
% TODO(Phase B): Figure figures/homeownership_by_age.jpg; the three prompts; two obstacles solved unprompted (FRED lacked the series -> Census HVS; HTTP 403 -> User-Agent); 6 min 21 s.
% TODO(Phase B): What the economist did: asked the question, checked the figure. Attribution: Paul Goldsmith-Pinkham, "From an Empty Folder to a Figure" + URL (Capstone_Projects_2026.md:526-528) + date.
\begin{frame}{Six Minutes, One Figure}
\textbf{Placeholder --- written in Phase B.} The specification is in the source comment above this frame.
\end{frame}

% ---- T1-F01 · EDIT · TIER: CORE · from T1:30
% TODO(Phase B): Name the economic twin of each part in the italic line (review A3).
\begin{frame}{Lecture 9 in One Map}
\textbf{One lecture, five parts, eleven decks, three questions: how agentic AI works, how to use it well, and what your role is.}

\LecNineMap{1}{}

\vspace{0.1cm}
\footnotesize\textit{Two research threads run through all five parts: \textbf{HA-Ramsey} (a published research program, full case in Lec10 T8) and the \textbf{ES Fellows dataset} (a project you could build this week, full case in Lec10 T4). Here we study the workflows; the cases themselves come later this morning.}
\end{frame}


%=============================================================================
\section{The Agent Opportunity}
%===========================================================

\sectiondivider{The Agent Opportunity}{What changed, and why an economist should care}

% ---- T1-F02 · EDIT · TIER: READ · from T1:45
\begin{frame}{Why Economists Should Care: Execution Is the Bottleneck}
\textbf{The bottleneck in applied research is often not ideas. It is execution.}

\vspace{0.3cm}

\begin{itemize}
    \item Turning a good question into clean code, diagnostics, tables, and figures is slow
    \item Many research tasks are repetitive: reshape data, rerun specifications, trace sources, refactor scripts, update plots
    \item Agentic tools compress the time from \textit{idea} to \textit{first defensible result}
    \item That raises the return to the skills economists uniquely provide: question design, institutional knowledge, validation, and judgment
\end{itemize}

\vspace{0.3cm}

\textbf{The promise is not ``AI does research for you.''} It is faster implementation and iteration around a question you still define and defend. The rest of the lecture supplies what that takes: how an agent runs (9.2), how to build and call one (9.3), how to design and validate one (9.4), and how to run, check, and govern a whole project (9.5).
\end{frame}

% ---- NEW · TIER: CORE
% TODO(Phase B): Review A4: one TikZ strip, six capability rungs, one course-owned example each, "what stayed human" underneath, and where each returns:
% TODO(Phase B): chat window (T1-F06 Fellows test) | one agent, one sitting (six minutes -> figure) | one agent, files, many sittings (882 Fellows) | agent + human gates, weeks (HA-Ramsey) | team of sub-agents (T3b ledger) | scripted fan-out + verifier (Lec10 T6b, 11,171 labels).
% TODO(Phase B): Counter-rung under "agent + gates": Lec10 T9, passed every gate, still wrong. Across all rungs: the GDP-by-party leading verb (T2a). Navier-Stokes is NOT on the ladder (dated sidebar in T3b and TA).
\begin{frame}{The Ladder: From a Chat Window to a Verified Swarm}
\textbf{Placeholder --- written in Phase B.} The specification is in the source comment above this frame.
\end{frame}

% ---- NEW · TIER: CORE
% TODO(Phase B): Review A3 lens frame. Researcher = principal, agent = agent: hidden action (Holmstrom 1979: you see the output, not the process) and a misspecified objective (Goodhart: the agent optimizes the contracted proxy).
% TODO(Phase B): "Generation is cheap; verification is scarce" (Tao, cited in Lec10 T9:693). One terminology line: three meanings of "agent" in this course (LLM in a loop; household in Aiyagari/HA-Ramsey; principal-agent theory).
% TODO(Phase B): Ownership: this frame owns the lens; Lec10 T9 ("This Is a Principal--Agent Problem", T9:607) owns the evidence -- point to it, do not retell it.
\begin{frame}{Delegating to an Agent Is a Principal--Agent Problem}
\textbf{Placeholder --- written in Phase B.} The specification is in the source comment above this frame.
\end{frame}

% ---- NEW · TIER: READ
% TODO(Phase B): Timeline TikZ (v1 T1 item 5): 2022 chat; 2023 structured tool calls + IDE assistants; 2024 terminal agents, MCP named once (T1b:71); 2025 sub-agents, frameworks at 1.0 (T1b:135-145), OpenAI adopts MCP (T1b:72), AAIF (T1b:73); 2026 scheduled runtimes (T1b:253-260).
% TODO(Phase B): Every fact [solid]/[hype]-tagged and dated; re-verify at the freshness gate. Closing line from T1b-F16.
\begin{frame}{From Chatbox to Agent Swarm, 2022--2026}
\textbf{Placeholder --- written in Phase B.} The specification is in the source comment above this frame.
\end{frame}

% ---- T1-F03 · KEEP · TIER: READ · from T1:62
\begin{frame}{Story 1 --- HA-Ramsey: Weeks, Not Years}
\textbf{Proper use of an agentic workflow turned one published solver into a new research program.}

\vspace{0.15cm}

\begin{examplebox}
\small
\textbf{Seed:} a working, validated representative-agent (RA) ``deep Ramsey'' neural solver --- the complete code behind Feng--Han--Zhu (2026).\\
\textbf{Target:} the heterogeneous-agent (HA) Ramsey model of Chien--Feng--Wen (2027) --- the entire paper: idiosyncratic risk, borrowing constraints, a 4--5D state.\\
\textbf{Outcome:} a full HA research codebase \emph{plus} three methodological discoveries that surfaced from the AI iteration itself --- $\alpha$-shapes for non-convex feasible sets, the admissible set as a viability kernel, and a Lyapunov-type penalty for saddle-path selection.
\end{examplebox}

\vspace{0.2cm}

\begin{itemize}
    \item The AI supplied \textbf{throughput} --- LaTeX drafting, PyTorch translation, diagnostics, 32+ sweeps; ``time from seed to working solver: \textbf{weeks, not years}; the binding constraint was economist judgment at the milestones, not implementation labor''
\end{itemize}

\textbf{Why show this first:} this is what ``using it properly'' buys. Not autocomplete --- research.

{\scriptsize \textit{Full case with the economics and the figures: Lec10 T8, later this morning. Source: HA-Ramsey case-study package.}}
\end{frame}

% ---- T1-F04 · EDIT · TIER: READ · from T1:85
\begin{frame}{Story 2 --- 882 Fellows: A Dataset You Could Build This Week}
\textbf{The second thread is deliberately ordinary: build a credible dataset about the economics profession from the public web.}

\vspace{0.25cm}

\begin{itemize}
    \item \textbf{Task:} career metadata for all 882 current Econometric Society Fellows --- election year, PhD year, research fields, career stage
    \item \textbf{Seed:} the official Society roster, split into 18 batches of 50; every filled field must carry a source URL from an official or institutional page
    \item \textbf{What makes it hard:} roster control, source discipline, state across sessions, auditable provenance --- \emph{bookkeeping, not intelligence}
    \item \textbf{What happened:} the agentic workflow scaled the search massively; the honest headline was about what still could \emph{not} be measured (T5c tells that part)
\end{itemize}

\vspace{0.25cm}

\textbf{Why two stories:} HA-Ramsey shows the ceiling of the method. The fellows dataset is the floor you can stand on immediately --- in Exercise B, you build one batch of 50 yourself.

\vspace{0.1cm}
{\scriptsize \textit{Full case with the coverage results: Lec10 T4, later this morning.}}
\end{frame}


%=============================================================================
\section{Chat, Copilot, Agent}
%===========================================================

\sectiondivider{Chat, Copilot, Agent}{Where the loop runs, and what it can touch}

% ---- T1-F05 · KEEP · TIER: READ · from T1:111
\begin{frame}{What Modern Web Chat Already Does Well}
\textbf{As of mid-2026, a strong web-chat product is already useful for economists.}

\vspace{0.25cm}

\begin{itemize}
    \item Brainstorming model setups, empirical strategies, review rubrics, and robustness checks
    \item Summarizing papers, referee reports, and lecture notes
    \item Inspecting uploaded files and explaining code snippets
    \item Sometimes browsing the web or running limited sandboxed tools
    \item Helping you draft a specification before you open your editor
\end{itemize}

\vspace{0.3cm}

\textbf{So the weak baseline is not ``a browser cannot do anything.''} The weak baseline is a \textit{single chat session without a durable harness}.
\end{frame}

% ---- T1-F06 · KEEP · TIER: READ · from T1:129
\begin{frame}{What Breaks Without a Harness: The Fellows Test}
\textbf{Try the fellows task in a plain chat window. It is a job of many sittings --- and four things break within the first hour.}

\vspace{0.2cm}

\begin{enumerate}
    \item \textbf{Durable context is weak.} The task spans days of separate sessions. The coding rules you settle today --- which sources count, how to record a missing PhD year --- are not reliably remembered by tomorrow's fresh chat
    \item \textbf{Reusable workflow is weak.} Your checklist lives in prompts you retype (slightly differently) each time, not in a file the tool reads automatically
    \item \textbf{Controlled execution is weak.} The chat cannot open the roster file, write the CSV, or rerun a check --- \emph{you} become the copy-paste machine between the model and your folder
    \item \textbf{Audit trail is weak.} Six weeks later: which source URL backed the PhD year of fellow \#312? A scrolled-away transcript cannot answer field-level questions
\end{enumerate}

\vspace{0.2cm}

\textbf{That is the real contrast.} Agents are not magic because they are ``smarter.'' They are useful because they operate inside a structured project with files, tools, and repeatable workflows.
\end{frame}

% ---- T1-F07 · MERGE-PENDING(T1-F07+T1-F10) · TIER: CORE · from T1:146
% TODO(Phase B): Review A5: merge into ONE frame keyed on where the loop runs, what it can touch, what persists, and who presses run (vendor sandbox vs IDE vs your machine/data/cluster).
% TODO(Phase B): Fix T1-F07's table: web chat runs code in a vendor sandbox (contradicted T1-F05); drop "why is Copilot not an agent?" -- IDE assistants ship agent modes; one product can sit on several rungs.
% TODO(Phase B): Delete T1-F07's JSON paragraph (T2a owns it). Name Claude Code, Codex CLI, Gemini CLI as terminal agents (review A2).
\begin{frame}{Chat vs.\ Agent: Same Brain, Different Arms}
\textbf{One key technical change made agents possible: modern LLMs (post-2023) support structured tool calls.}

\vspace{0.2cm}

Instead of writing ``you should run \texttt{python solve.py},'' the model emits a machine-readable request --- something like \texttt{\{"tool": "Bash", "command": "python solve.py"\}} --- and a \textbf{harness}, software installed on your machine, executes it and returns the real output. (T2a opens this loop in machine language.)

\vspace{0.25cm}

\begin{center}
\small
\begin{tabular}{|l|c|c|}
\hline
& \textbf{claude.ai (browser)} & \textbf{Claude Code (terminal)} \\ \hline
LLM model & \multicolumn{2}{c|}{the same frontier model} \\ \hline
Harness & none (sandboxed) & local CLI on your machine \\ \hline
Can run Python? & no & yes---\textit{your} Python, \textit{your} machine \\ \hline
Sees real errors? & no & yes---loops back automatically \\ \hline
\end{tabular}
\end{center}

\vspace{0.25cm}

\textbf{Same brain, different arms.} The terminal agent is not a magically smarter model. The harness is what flips the switch from text-only output to text\,+\,actions inside a project.
\end{frame}

% ---- T1-F10 · MERGE-PENDING(T1-F07+T1-F10) · TIER: CORE · from T1:228
\begin{frame}{The Spectrum, and Our Featured Tool}
\begin{center}
\small
\begin{tabular}{|p{2.5cm}|p{3.0cm}|p{3.2cm}|p{3.6cm}|}
\hline
& \textbf{Web chat} & \textbf{IDE assistant} & \textbf{Terminal agent} \\ \hline
\textbf{Best at} & thinking, drafting & editing local code & end-to-end execution \\ \hline
\textbf{Context} & conversation + uploads & open files & full project tree \\ \hline
\textbf{Typical role} & architect's whiteboard & coding copilot & workflow operator \\ \hline
\textbf{Economic example} & design the specification & clean one regression file & run the full pipeline \\ \hline
\end{tabular}
\end{center}

\vspace{0.25cm}

\textbf{Concrete example throughout this lecture:} Claude Code, a terminal agent. \textbf{General lesson:} the same workflow logic appears in Codex CLI, Gemini CLI, Aider, and similar tools.

\vspace{0.15cm}
\footnotesize\textit{The only tool-comparison table in Lecture 9 --- the workflow matters more than the tool. T5a returns to tool choice per pipeline step.}
\end{frame}

% ---- T1-F08 · KEEP · TIER: CORE · from T1:172
\begin{frame}{The Agentic Loop: Plan, Act, Observe, Verify, Repeat}
\textbf{An agentic workflow is not one long generation. It is a control loop wrapped around tools, files, and checks.}

\begin{center}
\begin{tikzpicture}[
    box/.style={rectangle, rounded corners, draw=RedTitle!70, fill=blue!5,
        minimum width=2.0cm, minimum height=0.75cm, text width=1.9cm,
        align=center, font=\scriptsize, text=DarkText},
    mem/.style={rectangle, rounded corners, draw=black!55, fill=VeryLightGray,
        minimum width=2.4cm, minimum height=0.8cm, text width=2.3cm,
        align=center, font=\scriptsize, text=DarkText},
    arr/.style={-{Stealth[length=2mm]}, thick, gray!60},
    lab/.style={font=\tiny\itshape, text=black!60}
]
    \node[box] (goal) at (-3.6,0.75) {User goal\\+ success criteria};
    \node[box] (plan) at (0,1.5)   {Plan\\next step};
    \node[box] (act)  at (3.1,0.75)  {Act\\shell, Python, files, web};
    \node[box] (obs)  at (3.1,-0.75) {Observe\\real outputs, real errors};
    \node[box] (ver)  at (0,-1.5)  {Verify\\tests, benchmarks, review};
    \node[mem] (memo) at (0,0)    {Project files\\+ memory on disk};
    \node[box] (art)  at (-3.6,-0.75) {Artifact\\+ report};
    \draw[arr] (goal) -- (plan);
    \draw[arr] (plan) to[bend left=18] (act);
    \draw[arr] (act) -- (obs);
    \draw[arr] (obs) to[bend left=18] (ver);
    \draw[arr] (ver) to[bend left=18] node[lab, below left] {not done yet: loop} (plan);
    \draw[arr] (ver) -- node[lab, below] {checks pass} (art);
    \draw[arr, dashed] (memo) -- (plan);
    \draw[arr, dashed] (act) -- (memo);
\end{tikzpicture}
\end{center}

\begin{itemize}
    \item \textbf{Orchestrator:} the loop deciding what to do next --- plan, act, observe, verify
    \item \textbf{Memory:} saved project state on disk, not just chat history (dashed arrows: the loop reads it and writes it)
    \item \textbf{Verifier:} tests, benchmarks, or review checks deciding whether output survives
\end{itemize}
\end{frame}

% ---- T1-F09 · KEEP · TIER: READ · from T1:212
% MiniCourse refinement: cleaner autonomy spectrum framing from M4_T3
\begin{frame}{What Makes Something ``Agentic''? Four Features}
\begin{itemize}
    \item \textbf{Tool use} --- invokes \texttt{pandas}, \texttt{curl}, \texttt{R}, \texttt{statsmodels}, the shell, whatever the task requires
    \medskip
    \item \textbf{Autonomy} --- plans a multi-step empirical workflow without re-prompting.\\
    \textit{e.g.,} ``download FRED $\to$ clean $\to$ regress $\to$ plot $\to$ save table'' as one command
    \medskip
    \item \textbf{Iteration} --- observes that a regression failed, identifies the bad column, fixes the formula, reruns
    \medskip
    \item \textbf{Persistence} --- remembers your project conventions across sessions (variable names, preferred packages, validated benchmarks)
\end{itemize}

\bigskip
\textbf{All four are cumulative.} Tool use without autonomy = a glorified interactive console (a REPL --- \emph{read--eval--print loop}, like the Python \texttt{>>>} prompt). Autonomy without iteration = a script that dies on the first error. You need all four for ``agentic.''
\end{frame}


%=============================================================================
\section{The Landscape in Three Frames}
%===========================================================

\sectiondivider{The Landscape in Three Frames}{Patterns, costs, and how long agents really work}

% ---- T1b-F04 · EDIT · TIER: READ · from T1b:102
\begin{frame}{Four Orchestration Patterns}
\textbf{Production agent systems are built from a small set of recurring control structures.}

\vspace{0.18cm}

\begin{center}
\footnotesize
\renewcommand{\arraystretch}{1.2}
\begin{tabular}{|p{3.0cm}|p{7.4cm}|}
\hline
\textbf{Pattern} & \textbf{Shape} \\ \hline
Orchestrator--worker & a lead agent decomposes a task and dispatches workers; \emph{dominant in production} \\ \hline
Planner--executor & one component plans, another executes step by step \\ \hline
Reflection & the agent critiques and revises its own output before returning \\ \hline
Handoff & control passes between specialized agents (e.g.\ triage $\to$ specialist) \\ \hline
\end{tabular}
\end{center}

\vspace{0.16cm}

\footnotesize\textbf{The agent team you will build in Part 9.3 and run as an RA team in T5b is an orchestrator--worker system}: the main agent orchestrates; \texttt{code-reviewer}, \texttt{domain-reviewer}, and \texttt{verifier} sub-agents are the workers. The landscape just adds vocabulary and tooling around the same loop.
\end{frame}

% ---- T1b-F06 · EDIT · TIER: READ · from T1b:152
% TODO(Phase B): Tag +90.2% / 15x as [hype] (vendor-internal eval) in the bullets themselves.
\begin{frame}{The Multi-Agent Cost--Benefit Decision}
\textbf{More agents is not free. Treat it as marginal value vs.\ marginal cost.}

\vspace{0.2cm}

\begin{columns}[T]
\begin{column}{0.56\textwidth}
\begin{itemize}
    \item Anthropic reports its multi-agent research system beat a single-agent baseline by \textbf{+90.2\%} on its internal research-eval\ldots
    \item \ldots at roughly \textbf{15$\times$ the token cost}
    \item So multi-agent pays off when the task is broad, parallelizable, and the answer's value is high relative to tokens
    \item It is wasteful for narrow, sequential, or cheap tasks
\end{itemize}

\vspace{0.12cm}
{\scriptsize \textbf{[solid, label: Anthropic-internal-eval]} --- a vendor benchmark on a self-defined research task, not an external suite. Read the $+90.2\%$ as directional, not portable.}
\end{column}
\begin{column}{0.42\textwidth}
\begin{takeawaybox}
\footnotesize\textbf{Economist framing.} You are choosing a point where marginal research value $=$ marginal token (and latency, and debugging) cost. The Aiyagari sub-agent review in T5b is justified on \emph{wall-clock}, not because ``more agents are better.''
\end{takeawaybox}
\end{column}
\end{columns}
\end{frame}

% ---- T1b-F09 · EDIT · TIER: READ · from T1b:217
% TODO(Phase B): Review A5: keep the concept (50%-reliability horizon, doubling time, fragility); move the model-specific numbers (Opus 4.5, GPT-5) to TA's dated comparisons and refresh them at the freshness gate.
% TODO(Phase B): Footer: "the rest of the landscape, dated: Appendix TA".
\begin{frame}{How Long Can an Agent Actually Work? Measure, Don't Trust Marketing}
\textbf{Vendors say ``works for hours'' or ``runs for days.'' METR gives a \emph{measured} alternative.}

\vspace{0.18cm}

\begin{columns}[T]
\begin{column}{0.56\textwidth}
\begin{itemize}\footnotesize
    \item METR's \emph{task-length-doubling} result: the task length an agent completes at \textbf{50\% reliability} has doubled roughly every \textbf{7 months}, accelerating to $\sim$\textbf{4 months} in 2024--25. \textbf{[solid]}
    \item Time Horizon 1.1 (Jan 2026) pegs \textbf{Opus 4.5 $\approx$ 320 min} and \textbf{GPT-5 $\approx$ 214 min} (50\%-reliability horizon).
    \item Use these measured minutes instead of vendor ``multi-hour/multi-day autonomy.''
\end{itemize}
\end{column}
\begin{column}{0.42\textwidth}
\begin{cautionbox}
\footnotesize\textbf{Metric fragility.} The number is tied to a \textbf{50\% reliability} threshold and a specific task suite; the horizon at 80\% reliability is much shorter, and it shifts across suites. A measured-but-fragile number still beats a vendor claim. \\[2pt]
\textbf{[hype counterpoint]} product pages advertising ``multi-hour'' or ``multi-day'' autonomy --- no reliability threshold stated.
\end{cautionbox}
\end{column}
\end{columns}

\vspace{0.12cm}
\footnotesize\textbf{Research translation:} the right question is not ``can it run for hours?'' but ``at what reliability, on what task distribution, validated how?''
\end{frame}


%=============================================================================
\section{What This Means for Research}
%===========================================================

\sectiondivider{What This Means for Research}{Four patterns, and what stays human}

% ---- T1-F11 · EDIT · TIER: READ · from T1:249
% TODO(Phase B): Review A3/A5: "Lec10's six cases" / "Cases 4--6" -> seven cases; add Case 7 (MPE shortcut, Lec10 T9) to the collaborator row as the cautionary instance.
\begin{frame}{Four Research Patterns Across the Lec10 Cases}
\small
\begin{tabular}{p{2.55cm}p{3.5cm}p{5.2cm}}
\toprule
\textbf{Pattern} & \textbf{Lec10 case(s)} & \textbf{Main risk if you use the wrong workflow} \\
\midrule
Reusable skill & Case 1: referee-style paper review & Generic comments, hallucinated quotes, missing quality gate \\
Empirical pipeline & Case 2: MEPS harmonization and figures & Wrong crosswalk, dropped survey weights, silent validation failure \\
Public-web dataset & Case 3: fellows metadata with field-level provenance & Weak sources, hidden missingness, no audit trail \\
Full-lifecycle collaborator & Cases 4--6: collaborator, 4-day sprint, HA-Ramsey & Polished output built on weak framing or claims discipline \\
\bottomrule
\end{tabular}

\vspace{0.05cm}
\textbf{Four patterns organize Lec10's six cases} --- and our two threads are instances: the fellows dataset (Case 3) and HA-Ramsey (Case 6, the full-lifecycle pattern at research scale).

\vspace{0.1cm}
\scriptsize\textbf{Where RAG fits.} Agentic RAG (Lec08 T5) differs from single-shot RAG: the agent decides \emph{dynamically}, via the same tool-call mechanism, whether and when to retrieve. RAG is one tool among many that the agent can compose.
\end{frame}

% ---- T1-F12 · MERGE-PENDING(T1-F12+T2-F24) · TIER: READ · from T1:269
% TODO(Phase B): Two rows of T2-F24 into T1-F12's columns; footer = one-line teaser of the seven tests (T5d owns the table); T5b's contract frame formally owns division of labor.
\begin{frame}{What Remains Human Across the Four Patterns}
\begin{columns}[T]
\begin{column}{0.48\textwidth}
\textbf{You still own}
\begin{itemize}
    \item the research question
    \item the success criteria
    \item the coding and source rules
    \item the validation target
    \item the final interpretation
\end{itemize}
\end{column}
\begin{column}{0.48\textwidth}
\textbf{The agent accelerates}
\begin{itemize}
    \item repetitive execution
    \item file handling and refactoring
    \item promptable decomposition
    \item reruns after fixes
    \item draft diagnostics and reports
\end{itemize}
\end{column}
\end{columns}

\vspace{0.3cm}
\textbf{Why this matters for economists:} the gain is not that judgment disappears. The gain is that disciplined execution becomes cheaper.

\vspace{0.15cm}
\textbf{Hold on to the left column.} In 9.5 it becomes a job description: the researcher as project manager.
\end{frame}

% ---- T2-F24 · MERGE-PENDING(T1-F12+T2-F24) · TIER: READ · from T2:783
\begin{frame}{Division of Labor: You Provide, the Agent Provides}
\begin{center}
\small
\begin{tabular}{|p{5.5cm}|p{5.9cm}|}
\hline
\textbf{You provide} & \textbf{Agent provides} \\ \hline
Question, identification logic, and source rules & Fast code generation and refactoring \\ \hline
Workflow choice and benchmark design & Boilerplate implementation and debugging \\ \hline
Knowledge of the data, model, and institutions & File edits, script writing, and test runs \\ \hline
Economic interpretation and seminar defense & Iteration on visible errors and outputs \\ \hline
\end{tabular}
\end{center}

\vspace{0.25cm}

\textbf{Research version of the slogan:} the agent can be a strong RA, but you still have to be the principal investigator.

\end{frame}

% ---- T1-F17 · EDIT · TIER: READ · from T1:436
% TODO(Phase B): Bridge sentence (v1): "You have seen what agents do and what proper use produced. The question changes from 'what is this?' to 'how does it actually run, and how do I start?'" + the 9.2 economic twin.
\begin{frame}{Bridge: Open the Hood}
\textbf{Part 9.1 framed the terminal-agent workflow: chat alone is not enough; a harness wraps tools, files, and verifiers around the model so the loop can do real work.}

\vspace{0.25cm}

\begin{itemize}
    \item You have seen \emph{what} agents do, \emph{why} the harness exists, and what proper use produced in two real projects
    \item \textbf{Dated detail (Appendix TA):} the wider landscape --- interoperability standards, frameworks, and product churn --- read on demand
    \item \textbf{Next (9.2, T2a and T2b):} the question changes from ``what is this?'' to ``\emph{how does it actually work, and how do I start?}'' --- components, install, initialization, the CLI, and your first session
\end{itemize}

\vspace{0.3cm}
\footnotesize\textit{Arc: \textbf{9.1 See} $\to$ 9.2 Under the Hood $\to$ 9.3 Build $\to$ 9.4 Design \& Validate $\to$ 9.5 Run \& Govern.}
\end{frame}

%===========================================================
\end{document}
